\documentclass[aspectratio=169, 11pt]{beamer}
\usepackage[utf8]{inputenc}
\usepackage[T1]{fontenc}
\usepackage[french]{babel}
\usepackage{graphicx}
\usepackage{tikz}
\usepackage{xcolor}
\usepackage{hyperref}

\usetheme{Madrid}
\usecolortheme{beaver}

\setbeamercolor{frametitle}{fg=black,bg=white}
\setbeamercolor{title in head/foot}{fg=black,bg=white}

\setbeamertemplate{footline}{%
  \leavevmode%
  \hbox{%
    \begin{beamercolorbox}[wd=.82\paperwidth,ht=2.25ex,dp=1ex,leftskip=1.2ex]{author in head/foot}%
      \usebeamerfont{author in head/foot}%
      \footnotesize RVP2 --- Moteur de recherche multimodal (Ping) --- Barbier / Ledoux / Pascal / Wone
    \end{beamercolorbox}%
    \begin{beamercolorbox}[wd=.18\paperwidth,ht=2.25ex,dp=1ex,right,rightskip=1.2ex]{page number in head/foot}%
      \usebeamerfont{page number in head/foot}%
      \footnotesize \insertframenumber/\inserttotalframenumber
    \end{beamercolorbox}%
  }%
}

\title[Highlights]{Fonctionnement du mode Highlights}
\author{Nathan Barbier - Projet Informatique}
\date{\today}

\begin{document}

\begin{frame}{Comment fonctionne le mode Highlights ?}
    \begin{columns}[T,onlytextwidth]
        \begin{column}{0.52\textwidth}
            {\small
            \begin{block}{Principe}
                Le backend appelle \texttt{/api/semantic/highlights}, filtre les points
                (match, gagnant, serveur, set), puis attribue a chaque point un
                \textbf{score explicite sur 100}.
            \end{block}

            \begin{block}{Score calcule}
                \textbf{Highlight score} =
                \[
                \sum_i w_i \times s_i
                \]
                avec des poids normalises. Par defaut :
                \textbf{duree 35\%}, \textbf{profondeur du rallye 15\%},
                \textbf{variete des effets 15\%},
                \textbf{variete des lateralites 10\%},
                \textbf{variete des zones 10\%},
                \textbf{fin du point 10\%},
                \textbf{pression du score 5\%}.
            \end{block}

            \begin{block}{Ce que le score regarde concretement}
                \begin{itemize}
                    \item \textbf{Duree / nb de coups} : un echange long monte dans le classement.
                    \item \textbf{Variete} : plus la sequence change d'effets, de zones ou de lateralites, plus le point est riche.
                    \item \textbf{Fin du point} : \texttt{pt\_gagne} est tres valorise ; un dernier coup agressif (\texttt{topspin}, \texttt{flip}) ajoute encore du score.
                    \item \textbf{Contexte} : un point serre en fin de set, ou une balle de set, augmente la note.
                \end{itemize}
            \end{block}

            \begin{alertblock}{Resultat}
                Les points sont tries par \texttt{highlight\_score} decroissant :
                le mode Highlights ne cherche donc pas un texte similaire,
                il \textbf{classe automatiquement les plus beaux points}.
            \end{alertblock}
            }
        \end{column}

        \begin{column}{0.46\textwidth}
            \vspace{0.2cm}
            \begin{beamercolorbox}[wd=\linewidth,ht=0.68\textheight,dp=0.2cm,center,rounded=true,shadow=true]{block body}
                \vfill
                {\Large Screenshot du mode Highlights}

                \vspace{0.3cm}
                {\normalsize A inserer ici}
                \vfill
            \end{beamercolorbox}

            \vspace{0.2cm}
            {\scriptsize
            Astuce orale : expliquer que les ponderations sont reglables dans l'interface,
            puis montrer visuellement comment le classement change.}
        \end{column}
    \end{columns}
\end{frame}

\end{document}
